%@ pages 1-3
%@ status translated
% Title page, abstract, translator's note and contents.
% The original's title, authors, English abstract and French "Résumé"
% occupy pp. 1--3; the hand-written Sommaire follows on p. 3.
%
% The English abstract is the authors' own text and is reproduced verbatim.
% The French "Résumé" is not reprinted: it duplicates the abstract.

\maketitle

\noindent
\textit{Originally published in French as ``Descente pour les $n$-champs'',
arXiv:math/9807049v3 (13 March 2001). This is an English translation of the
body; the French original remains authoritative.}

\bigskip

\begin{abstract}
\noindent
We develop the theory of $n$-stacks (or more generally Segal $n$-stacks
which are $\infty$-stacks such that the morphisms are invertible above
degree $n$). This is done by systematically using the theory of
closed model categories (cmc). Our main results are: a definition of
$n$-stacks in terms of limits, which should be perfectly general for stacks
of any type of objects; several other characterizations of $n$-stacks in
terms of ``effectivity of descent data''; construction of the stack
associated to an $n$-prestack; a strictification result saying that any
``weak'' $n$-stack is equivalent to a (strict) $n$-stack; and a descent
result saying that the $(n+1)$-prestack of $n$-stacks (on a site) is an
$(n+1)$-stack. As for other examples, we start from a ``left Quillen
presheaf'' of cmc's and introduce the associated Segal $1$-prestack. For
this situation, we prove a general descent result, giving sufficient
conditions for this prestack to be a stack. This applies to the case of
complexes, saying how complexes of sheaves of $\Oo$-modules can be glued
together via quasi-isomorphisms. This was the problem that originally
motivated us.
\end{abstract}

\bigskip

\noindent
\textit{Translator's note.} The authors' v3 errata header, which records an
error in Proposition~6.1/Lemma~6.2 and its consequences for 10.11, 12.1,
14.4, 15.8, 19.4 and 21.1, is preserved in the source and is reproduced in
the relevant places. The mathematical content itself has not been
independently checked. French terminology is translated as recorded in
\texttt{terminology.tex}; a few terms are kept in French at first
occurrence because they are part of the history.

\newpage

\begin{center}
{\Large \underline{Contents}}
\end{center}

\noindent
{\bf 1. Introduction}---p. \pageref{intropage}
\newline
{\small
Motivation, description of the results.
}

\noindent
{\bf 2. Segal $n$-categories}---p. \pageref{ncatsegpage}
\newline
{\small
Definition of Segal $n$-categories, truncation, interior. The general
method for obtaining a cmc structure, and its application to Segal
$n$-categories. Homotopy classes of morphisms.}

\noindent
{\bf 3. Segal $n$-prestacks}---p. \pageref{nchampsegpage}
\newline
{\small
The cmc structure for Segal $n$-prestacks (notions of weak equivalence,
cofibration, fibration), comparison with the coarse topology, compatibility
with truncation, properness. Descent data.}

\noindent
{\bf 4. The functorialities $p^{\ast}$, $p_{\ast}$, $p_!$}---p.
\pageref{fonctorialitepage}
\newline
{\small
Relation between Segal $n$-prestacks on $\Xx$ and $\Zz$ for a functor
$p:\Zz\rightarrow \Xx$, criteria for preservation of fibrant objects.}

\noindent
{\bf 5. The Bousfield--Kan-type structure after Hirschhorn}---p.
\pageref{bousfieldkanpage}
\newline
{\small
The HBKQ cmc structure for Segal $n$-prestacks, in which the cofibrations
are obtained by freely adding cells; application to the notion of descent
data.}

\noindent
{\bf 6. The localisation viewpoint}---p.
\pageref{localisationpage}
\newline
{\small
The cmc's of Segal $n$-prestacks for the topology $\Gg$ can be obtained
from the cmc's for the coarse topology by inverting the arrows corresponding
to the sieves of $\Gg$.}

\noindent
{\bf 7. Segal categories and simplicial categories}---p.
\pageref{segsimplpage}
\newline
{\small
Relation between Segal $1$-categories and simplicial categories; homotopy
products and coproducts and the Gabriel--Zisman argument.}

\noindent
{\bf 8. Dwyer--Kan localisation}---p.
\pageref{dwyerkanpage}
\newline
{\small
Dwyer--Kan localisation provides many examples of simplicial categories,
and moreover makes it possible to show that every Segal $1$-category is
equivalent to a simplicial category.}

\noindent
{\bf 9. A first definition of stack}---p.
\pageref{notionchamppage}
\newline
{\small
Definition of Segal $n$-stack as a Segal $n$-prestack $A$ whose $A(X)$ are
Segal $n$-categories and for which the morphism $A\rightarrow A'$ to the
$\Gg$-fibrant replacement is an object-by-object equivalence. Notion of
associated stack, compatibilities.}

\noindent
{\bf 10. Criteria for a prestack to be a stack}---p.
\pageref{caracterisationpage}
\newline
{\small
$A$ is a Segal $n$-stack if and only if the $A_{1/}(x,y)$ are Segal
$n-1$-stacks on $\Xx /X$ for all $x,y\in A_0(X)$, and if the descent data
for a sieve $\Bb \subset \Xx /X$ are effective. Several versions of this
criterion are given.}

\noindent
{\bf 11. Internal model categories}---p.
\pageref{cmipage}
\newline
{\small
We define the notion of internal model category and use it to define the
$n+1$-prestack $nSe\underline{CHAMP}(\Xx )$ of Segal $n$-stacks over $\Xx$,
as well as its Segal $n+1$-category $nSeCHAMP(\Xx )$ of global sections.
Comparison with the Dwyer--Kan localisation.}

\noindent
{\bf 12. The universal family}---p.
\pageref{universellepage}
\newline
{\small
We define a morphism $nSeCHAMP(\Xx )\rightarrow \underline{Hom}(\Xx ^o,
nSeCAT')$, show that it is fully faithful (and state the result
\ref{correlation}, which identifies its essential image). This leads us to
a new version of the notion of stack.}

\noindent
{\bf 13. The stack associated to a prestack}---p.
\pageref{associepage}
\newline
{\small
We define the functor ``stack associated to a prestack'' from
$nSeCHAMP(\Xx ^{\rm gro})$ to $nSeCHAMP (\Xx ^{\Gg})$, which is adjoint to
the inclusion functor in the other direction.}

\noindent
{\bf 14. Limits}---p.
\pageref{limitespage}
\newline
{\small
Computation of limits using cmc structures; definition of stacks in terms
of limits.}

\noindent
{\bf 15. A little more on the descent condition}---p.
\pageref{peupluspage}
\newline
{\small
We redo the criteria of \S10 in terms of limits, and make explicit the
computation of these limits in the case of a sieve defined by a covering
family, in terms of limits taken over $\Delta$.}

\noindent
{\bf 16. The Grothendieck construction}---p.
\pageref{gropage}
\newline
{\small
We construct the ``$n$-fibred category'' associated to an $n$-prestack over
a category, an operation denoted like an integral following Thomason. We
compare the sections of the $n$-prestack and those of the associated
$n$-fibred category.}

\noindent
{\bf 17. Quillen presheaves}---p.
\pageref{quipage}
\newline
{\small
We introduce the notion of left Quillen presheaf, a kind of presheaf of
cmc's in which the restriction functors are left Quillen functors. We give
various cmc structures for the sections of the integral of a Quillen
presheaf.}

\noindent
{\bf 18. Strictification}---p.
\pageref{strictpage}
\newline
{\small
We strictify the weak sections of Segal $1$-prestacks which come from
Quillen presheaves by Dwyer--Kan localisation. This makes it possible to
strictify weak families of Segal $n$-categories into Segal $n$-prestacks.
End of the proof of \ref{correlation}.}

\noindent
{\bf 19. Descent for left Quillen presheaves}---p.
\pageref{quidescentepage}
\newline
{\small
For a left Quillen presheaf $\mm$ we give a criterion for the associated
Segal $1$-prestack $L(\mm )$ to be a stack.}

\noindent
{\bf 20. Example: descent for $n$-stacks}---p.
\pageref{exemplepage}
\newline
{\small
We use the criterion of \S19 to prove that $nSe\underline{CHAMP}(\Xx )$ is a
Segal $n+1$-stack. This generalises the classical theorem on the gluing of
sheaves. We also give a direct proof.}

\noindent
{\bf 21. Example: descent for complexes}---p.
\pageref{complexepage}
\newline
{\small
We use the criterion of \S19 to prove that $L(Cpx _{\Oo},qis)$ is a Segal
$1$-stack, for a reasonable ringed site $(\Xx , \Oo )$. We state a
geometricity result for the moduli stack of perfect complexes.}

\bigskip

\noindent
\textit{References}---p. \pageref{references}.
